% ======================================================================
% SECTION 6: AUTONOMOUS REGULATORY AFFAIRS AND SUBMISSIONS
% File: sections/regulatory_affairs.tex
% ======================================================================

Regulatory affairs automation replaces the traditional regulatory team with a dedicated
agent that manages the IND lifecycle, automates multi-jurisdiction filing, navigates
accelerated review pathways, and monitors the evolving regulatory landscape. The
regulatory agent operates within the adapted frameworks established by prior work in this
repository: 21 CFR Part 312 adapted for Physical AI \cite{cfr312-adapt}, 21 CFR Part 50
adapted for human subjects protection in robotic procedures \cite{cfr50-adapt}, and ICH
E6(R3) adapted for quality-by-design in Physical AI trials \cite{ich-adapt}.

\subsection{Regulatory Agent}

The regulatory agent (regulatory\_agent) replaces the regulatory affairs team for IND
lifecycle management, authority correspondence, and multi-jurisdiction coordination. Under
21 CFR \S312.50, the sponsor retains full responsibility for regulatory interactions
regardless of which activities are delegated to CROs or other agents
\cite{cfr-part-312}. The regulatory agent embodies this non-delegable accountability
through direct management of all regulatory submissions and authority communications.

IND lifecycle management spans the complete regulatory lifecycle from pre-IND planning
through post-marketing commitments. Initial IND filing follows 21 CFR Part 312 Subpart B
requirements: the agent assembles the IND application from protocol, Investigator's
Brochure, CMC data, nonclinical study summaries, and clinical investigator information
generated by the respective specialized agents. The agent tracks the 30-day FDA review
window and manages the safe-to-proceed determination. Subsequent IND amendments per
\S312.30 are generated automatically when protocol changes, new safety information, or
CMC updates warrant notification. Annual reports per \S312.33 are compiled from
cumulative study data with automated generation of required tables and narratives.

Authority correspondence manages the formal meeting process with regulatory agencies.
For FDA interactions, the agent prepares meeting request packages for Type A (immediately
needed), Type B (standard milestone), and Type C (other) meetings. Meeting briefing
documents are generated with structured sections covering the meeting objectives,
background, questions for discussion, and proposed approaches. Post-meeting minutes are
tracked against meeting agreements with automated implementation of regulatory
commitments.

Multi-jurisdiction coordination manages simultaneous regulatory processes across the FDA,
European Medicines Agency, Pharmaceuticals and Medical Devices Agency (Japan), Health
Canada, Therapeutic Goods Administration (Australia), and other relevant authorities. The
agent maintains jurisdiction-specific regulatory intelligence, tracks divergent requirements,
and generates jurisdiction-appropriate documents from a common source dataset.
Table~\ref{tab:regulatory-filing} presents the global regulatory filing matrix.

\begin{longtable}{L{1.5cm} L{1.8cm} L{1.5cm} L{3.2cm} L{3.2cm}}
\caption{Global Regulatory Filing Matrix}
\label{tab:regulatory-filing} \\
\toprule
\textbf{Jurisdiction} & \textbf{Filing Type} & \textbf{Timeline} & \textbf{Key Requirements} & \textbf{Agent Automation} \\
\midrule
\endfirsthead
\toprule
\textbf{Jurisdiction} & \textbf{Filing Type} & \textbf{Timeline} & \textbf{Key Requirements} & \textbf{Agent Automation} \\
\midrule
\endhead
\midrule
\multicolumn{5}{r}{\textit{Continued on next page}} \\
\endfoot
\bottomrule
\endlastfoot
United States (FDA) & IND & 30-day review & 21 CFR 312, eCTD v4.0, clinical hold response within 30 days & Automated eCTD assembly, serial submission tracking, amendment generation \\
European Union (EMA) & CTA via CTIS & Variable by member state & EU CTR 536/2014, CTIS structured submission, member state assessment & Automated CTIS data entry, parallel member state tracking, lay summary generation \\
Japan (PMDA) & CTN & 30-day notification & PMDA CTN format, Japanese-language requirements & Template-based CTN generation, PMDA meeting preparation \\
Canada (HC) & CTA & 30-day default & Health Canada CTA format, NOL response, clinical hold provisions & Automated CTA package assembly, NOL tracking \\
Australia (TGA) & CTN/CTA & CTN immediate, CTA 30-day & TGA notification or approval based on risk & Risk-based pathway selection, automated notification \\
China (NMPA) & IND & 60-day review & NMPA IND requirements, local clinical data requirements & Parallel filing with FDA, China-specific CMC supplements \\
\end{longtable}

\subsection{IND and CTA Filing Automation}

Filing automation addresses the document assembly and submission process across jurisdictions.
The eCTD v4.0 standard (effective September 2024) \cite{fda-ectd} provides the technical
framework for electronic regulatory submissions. The agent assembles eCTD modules from
component documents generated by specialized agents: Module 1 (administrative and
prescribing information) from the regulatory agent, Module 2 (summaries) from the writing
agent, Module 3 (quality/CMC) from the supply agent, Module 4 (nonclinical) from the
clinical accountability agent, and Module 5 (clinical) from the data and biostatistics
agent.

The EU Clinical Trials Information System \cite{eu-ctis} transition (mandatory as of 31
January 2025 for all ongoing trials) requires structured electronic submission through the
CTIS portal rather than national authority submission. The regulatory agent manages the CTIS
workflow including: initial application submission with Part I (common assessment) and Part
II (member state-specific), substantial modification submissions, and results posting (within
one year of trial completion, six months for pediatric trials). The EU Clinical Trials
Regulation \cite{eu-ctr} 25-year archiving requirement is managed through integration with
the archive manager (\S\ref{sec:writing}).

UK CTIMP reforms taking effect in April 2026 affect documentation and contracting
requirements for clinical trials of investigational medicinal products. The regulatory agent
maintains UK-specific regulatory intelligence and generates UK-compliant documents when
UK sites are included in the trial.

The eCTD lifecycle management extends beyond initial filing to encompass the ongoing
maintenance of the regulatory dossier throughout the trial lifecycle. Each protocol
amendment triggers an eCTD sequence update. Each new safety report generates an IND
safety report sequence. Manufacturing changes produce CMC supplement sequences. The
regulatory agent maintains the complete eCTD sequence history, ensuring that the
cumulative dossier accurately reflects the current state of the program at any point in
time. This eCTD lifecycle management is essential for RTOR (Real-Time Oncology Review),
which requires the submission of individual eCTD modules as they become available rather
than waiting for the complete application.

The EU CTIS requirements represent a particularly complex automation challenge because the
system requires structured data entry (not document upload) for many application elements.
The regulatory agent maps the protocol content, investigator information, and site details
to CTIS data fields, validates the data against CTIS business rules before submission, and
manages the parallel assessment process across member states. The Part I (common assessment)
and Part II (member state-specific assessment) workflows are tracked independently, with
cross-referencing to ensure consistency between the common and member state-specific
elements.

CDISC standards compliance \cite{cdisc-sdtm, cdisc-adam} ensures that all submission
datasets meet FDA standardized study data requirements. The agent validates SDTM domains
and ADaM datasets against current implementation guides and generates Define.xml metadata
files. Pinnacle 21 (P21) validation checks are executed automatically before any dataset
is included in a regulatory submission package.

\subsection{Accelerated Pathway Navigation}

The accelerated pathway navigation subsystem evaluates program eligibility for expedited
review mechanisms across jurisdictions and generates designation request packages when
criteria are met. Table~\ref{tab:accelerated-pathways} compares the major accelerated
pathways available for oncology programs.

\begin{longtable}{L{2.0cm} L{1.5cm} L{3.0cm} L{2.2cm} L{3.0cm}}
\caption{Accelerated Pathway Comparison}
\label{tab:accelerated-pathways} \\
\toprule
\textbf{Pathway} & \textbf{Jurisdiction} & \textbf{Eligibility} & \textbf{Benefit} & \textbf{Agent Decision Criteria} \\
\midrule
\endfirsthead
\toprule
\textbf{Pathway} & \textbf{Jurisdiction} & \textbf{Eligibility} & \textbf{Benefit} & \textbf{Agent Decision Criteria} \\
\midrule
\endhead
\midrule
\multicolumn{5}{r}{\textit{Continued on next page}} \\
\endfoot
\bottomrule
\endlastfoot
Breakthrough Therapy & FDA & Substantial improvement over existing therapies for serious conditions & Intensive FDA guidance, rolling review eligibility & Preliminary clinical evidence of substantial improvement over available therapy \\
Fast Track & FDA & Serious condition with unmet medical need & More frequent FDA meetings, rolling review & Evidence of unmet need and potential to address it \\
Accelerated Approval & FDA & Surrogate endpoint reasonably likely to predict clinical benefit & Earlier approval on surrogate endpoint & Validated or reasonably likely surrogate endpoint available \\
Priority Review & FDA & Significant improvement in safety or effectiveness & 6-month review (vs. standard 10-month) & Evidence of significant improvement based on submitted data \\
Project Orbis & FDA-led & Oncology products with international interest & Concurrent review by multiple agencies & Multi-national filing readiness, participating authority agreement \\
RTOR & FDA & Oncology products with potential for significant benefit & Real-time iterative review of submission modules & Submission modules ready for sequential review \\
PRIME & EMA & Substantial clinical improvement for unmet medical need & Enhanced EMA scientific advice, accelerated assessment eligibility & Compelling non-clinical or early clinical data \\
Conditional MA & EMA & Unmet medical need, positive benefit-risk on comprehensive data & Renewable 1-year marketing authorization & Positive benefit-risk with obligation to provide confirmatory data \\
SAKIGAKE & PMDA & Innovative products developed in Japan & Prioritized review, pre-application consultation & Japan-first or Japan-simultaneous development \\
NOC/c & Health Canada & Serious condition, surrogate or clinical endpoint & Conditional approval with post-market obligations & Evidence supporting conditional approval \\
\end{longtable}

The pathway evaluation algorithm scores each active program against the eligibility criteria
for all available pathways. When a program meets the criteria for one or more pathways, the
agent generates a recommendation ranking that considers the benefit magnitude, application
burden, and strategic alignment with the overall development plan. Designation request
packages are pre-assembled so that applications can be submitted quickly when the portfolio
agent approves the pathway strategy.

Project Orbis \cite{fda-project-optimus} concurrent multi-country review and Real-Time
Oncology Review enable accelerated international access for promising oncology therapies.
The regulatory agent manages the Orbis process by coordinating submission timing across
participating authorities and tracking parallel review timelines.

\subsection{Regulatory Intelligence and Horizon Scanning}

The regulatory intelligence subsystem monitors regulatory guidance repositories across
jurisdictions for updates that may affect active programs. Key monitoring targets include
FDA draft and final guidance documents, ICH guideline revisions, EMA scientific advice
updates, and jurisdiction-specific regulatory changes.

Impact assessment algorithms evaluate each detected regulatory change against the portfolio
of active programs. For each affected program, the agent generates an impact assessment
documenting the regulatory change, the affected program elements, the compliance gap (if
any), the recommended response, and the implementation timeline. High-impact changes (those
affecting primary endpoints, safety reporting requirements, or submission timelines) are
escalated immediately. Low-impact changes enter the standard regulatory intelligence queue
for periodic review.

The regulatory intelligence subsystem maintains a knowledge graph of regulatory requirements
that represents the relationships between regulations, guidance documents, and operational
requirements. When a new guidance document is detected, the system identifies which nodes in
the knowledge graph are affected and propagates the impact assessment to all connected
operational requirements. This structured approach ensures that regulatory changes are not
evaluated in isolation but are assessed in the context of the complete regulatory environment.

For Physical AI trials specifically, the regulatory intelligence subsystem tracks the
evolving regulatory landscape for AI/ML in clinical trials, robotic surgical devices, digital
health technologies, and federated data sharing. These regulatory domains are developing
rapidly, and changes in any one domain may affect the autonomous sponsor's operations. The
agent maintains separate tracking channels for FDA, EMA, PMDA, Health Canada, and TGA
positions on these topics, enabling the system to anticipate regulatory convergence or
divergence across jurisdictions.

The regulatory agent also manages the interface between the investigational drug regulatory
pathway and the device regulatory pathway. Physical AI trials involve both investigational
drugs (regulated under 21 CFR Part 312 \cite{cfr-part-312}) and medical devices (regulated
under 21 CFR Parts 510(k), De Novo, or PMA pathways). When a trial involves a
combination product (drug plus device), the regulatory agent coordinates with both CDER
(or CBER) and CDRH divisions of the FDA, managing the dual submission pathway and ensuring
that device-specific requirements (biocompatibility testing, electromagnetic compatibility,
software validation) are integrated with the drug development program.

Compliance calendar management tracks all regulatory deadlines across programs and
jurisdictions. The calendar integrates IND annual report deadlines, DSUR deadlines, CTIS
submission windows, E2B(R3) transmission deadlines, ClinicalTrials.gov results posting
deadlines (one year from primary completion per FDAAA 801 \cite{fdaaa-801}), and all other
jurisdiction-specific deadlines. Automated reminders are generated at configurable lead
times, and deadline compliance is tracked as a KPI reported to the quality agent.
